\documentclass[11pt]{article}

\usepackage[margin=1in]{geometry}
\usepackage{amsmath,amssymb}
\usepackage{booktabs}
\usepackage{enumitem}
\usepackage{hyperref}
\hypersetup{hidelinks}

\title{Scaling Book Exercises: Section 12, Quiz 1}
\author{Lucas Yan}
\date{}

\begin{document}
\maketitle

\noindent
\textbf{Quiz:} GPU Hardware\\
\textbf{Source scan:} \texttt{../handwritten/section\_12\_handwritten.pdf},
pages 1--3\\
\textbf{Reference:} \url{https://jax-ml.github.io/scaling-book/gpus/}

\section*{Question 1: CUDA cores}

An H100 has 132 streaming multiprocessors (SMs), four subpartitions per SM,
and 32 fp32 CUDA cores per subpartition:
\[
\boxed{132\cdot4\cdot32=16{,}896\text{ CUDA cores}.}
\]
A B200 has 148 SMs:
\[
\boxed{148\cdot4\cdot32=18{,}944\text{ CUDA cores}.}
\]
A TPU v5p has two TensorCores, four independent ALUs per VPU lane, and
\(8\cdot128\) lanes per TensorCore:
\[
\boxed{2\cdot4\cdot8\cdot128=8{,}192\text{ ALUs}.}
\]
The TPU therefore has roughly half as many independent vector ALUs as an
H100.

\section*{Question 2: H100 vector FLOPs}

Ignoring fused multiply-add counting, the base-clock vector throughput is
\[
132\cdot4\cdot32\cdot1.59\times10^9
=\boxed{26.9\ \text{TFLOPs/s}}.
\]
At the \(1.98\) GHz boost clock it is
\[
132\cdot4\cdot32\cdot1.98\times10^9
=\boxed{33.5\ \text{TFLOPs/s}}.
\]
The H100 has \(990\) bf16 Tensor-Core TFLOPs/s, about \(30\times\) the
non-FMA vector rate.

\section*{Question 3: Matmul intensity}

The critical arithmetic intensity is peak matmul throughput divided by HBM
bandwidth:
\begin{align*}
I_{\mathrm{H100,fp16}}
&=\frac{990\times10^{12}}{3.35\times10^{12}}
=\boxed{295\ \text{FLOPs/byte}},\\
I_{\mathrm{B200,fp16}}
&=\frac{2250\times10^{12}}{8\times10^{12}}
=\boxed{281\ \text{FLOPs/byte}}.
\end{align*}
FP8 doubles the matmul throughput, giving approximately
\[
\boxed{590\text{ FLOPs/byte on H100},\qquad
562\text{ FLOPs/byte on B200}.}
\]

\section*{Question 4: B200 matmul runtimes}

For
\(\mathrm{fp16}[64,4096]\times\mathrm{fp16}[4096,8192]\),
the effective batch \(64<281\), so the operation is bandwidth-bound.  It
moves
\[
2BD+2DF+2BF
=2(64)(4096)+2(4096)(8192)+2(64)(8192)
\approx69\times10^6\text{ bytes}.
\]
The ideal runtime is therefore
\[
\boxed{T\approx\frac{69\times10^6}{8\times10^{12}}
=8.6\ \mu\mathrm{s}},
\]
or roughly \(10\)--\(12\,\mu\mathrm{s}\) in practice.

For batch \(512>281\), the matmul is compute-bound:
\[
\boxed{
T\approx
\frac{2(512)(4096)(8192)}{2.25\times10^{15}}
\approx15.3\ \mu\mathrm{s}},
\]
with roughly \(20\,\mu\mathrm{s}\) being a practical expectation.

\section*{Question 5: H100 fast-memory capacity}

Each of the 132 SMs has \(256\) kB of SMEM and \(256\) kB of register
memory:
\[
132(256\text{ kB})\approx\boxed{33\text{ MB of SMEM}}
\]
and independently about \(33\) MB of registers, or \(66\) MB combined.
A modern TPU has about \(120\) MB of VMEM but only \(256\) kB of vector
registers in total.

\section*{Question 6: B200 clock frequency}

The \(18{,}944\) CUDA cores can perform two FLOPs per cycle with an FMA:
\[
18{,}944\cdot2=37{,}888\ \text{FLOPs/cycle}.
\]
Thus the clock implied by \(80\) vector fp32 TFLOPs/s is
\[
\boxed{
f=\frac{80\times10^{12}}{37{,}888}
\approx2.1\ \text{GHz}.}
\]

\section*{Question 7: H100 vector-add runtime}

Adding two fp32 vectors of length \(N\) performs \(N\) FLOPs.  Two inputs
are read and one output is written, moving
\[
4N+4N+4N=12N\text{ bytes}.
\]
Hence
\[
\boxed{I=\frac{N}{12N}=\frac1{12}\ \text{FLOP/byte}.}
\]
The two roofline times are
\[
T_{\mathrm{math}}=\frac{N}{33.5\times10^{12}},
\qquad
T_{\mathrm{comms}}=\frac{12N}{3.35\times10^{12}}
=\frac{N}{2.79\times10^{11}}.
\]
Because \(T_{\mathrm{comms}}\gg T_{\mathrm{math}}\), vector addition is
strongly memory-bound.  For \(N=1024\), the ideal full-device roofline gives
\[
T_{\mathrm{math}}\approx0.031\text{ ns},
\qquad
T_{\mathrm{comms}}\approx3.67\text{ ns}.
\]
Only 1024 CUDA cores, or about eight SMs, can be occupied at this size, so
the compute estimate would be roughly \(16\times\) larger even before kernel
launch latency.  In practice, fixed latency dominates both estimates.

For \(N=1024^3\),
\[
T_{\mathrm{math}}\approx32.0\ \mu\mathrm{s},
\qquad
\boxed{T_{\mathrm{roofline}}\approx T_{\mathrm{comms}}
\approx3.85\text{ ms}.}
\]
No empirical GPU measurement is claimed in the handwritten work.

\end{document}
